Este trabalho apresenta o desenvolvimento de um sistema fuzzy Takagi-Sugeno (TS) de ordem zero e primeira ordem para a aproximação da função não linear $f(x) = e^{-x/5} \cdot \sin(3x) + 0.5 \cdot \sin(x)$ no intervalo $x \in [0, 10]$. Foram testadas diferentes funções de pertinência (gaussianas, triangulares e trapezoidais) e métodos de ajuste dos consequentes (média ponderada, gradiente descendente e Recursive Least Squares - RLS). O desempenho do sistema foi avaliado utilizando o RMSE, e os resultados mostram que a abordagem TS de primeira ordem com funções de pertinência gaussianas apresentou o menor erro de aproximação. O ajuste dos consequentes por gradiente descendente e RLS também contribuiu para a redução do erro na ordem zero. Os experimentos demonstram a eficácia do sistema fuzzy TS na modelagem de funções não lineares complexas.